%!TEX root = ../../main.tex
% ============================================================
% 数据表格 / 购买与运算表格
% 本文件由 main.tex 通过 \input 引入，请勿单独编译
% 重构日期: 2026-04-11
% ============================================================

\subsection{解答：世界上陆地面积最大的四个国家}

\vspace{0.4cm}

\noindent\textbf{4. 世界上陆地面积最大的四个国家的面积如下表：}

\vspace{0.4cm}

\begin{center}
\renewcommand{\arraystretch}{1.6}
\setlength{\tabcolsep}{15pt}
\begin{tabular}{|c|c|}
\hline
中\quad 国 & 九百六十万平方千米 \\
\hline
俄罗斯 & 一千七百零九万八千二百平方千米 \\
\hline
美\quad 国 & 九百三十七万平方千米 \\
\hline
加拿大 & 九百九十八万平方千米 \\
\hline
\end{tabular}
\end{center}

\vspace{0.6cm}

\noindent\textbf{（1）填写下表。}

\vspace{0.4cm}

\begin{center}
\renewcommand{\arraystretch}{1.6}
\setlength{\tabcolsep}{12pt}
\begin{tabular}{|c|c|c|}
\hline
国\quad 家 & 面积/平方千米 & 面积/万平方千米 \\
\hline
中\quad 国 & \textcolor{red}{9600000} & \textcolor{red}{960} \\
\hline
俄罗斯 & \textcolor{red}{17098200} & \textcolor{red}{1709.82} \\
\hline
美\quad 国 & \textcolor{red}{9370000} & \textcolor{red}{937} \\
\hline
加拿大 & \textcolor{red}{9980000} & \textcolor{red}{998} \\
\hline
\end{tabular}
\end{center}

\vspace{0.6cm}
{\small\color{gray}
\textbf{解析：}\\
首先写出各面积的阿拉伯数字形式：\\
九百六十万 $\rightarrow$ 9600000\\
一千七百零九万八千二百 $\rightarrow$ 17098200\\
九百三十七万 $\rightarrow$ 9370000\\
九百九十八万 $\rightarrow$ 9980000\\
将其改写为以“万”为单位的数，只需将小数点向左移动 4 位：\\
$9600000 = 960$ 万；
$17098200 = 1709.82$ 万；
$9370000 = 937$ 万；
$9980000 = 998$ 万。
}

\vspace{0.8cm}

%% ==============================
%% 分数的初步认识：涂色表示分数
%% ==============================
\newpage

\newpage

\subsection{解答：体育用品商店购买问题}

\vspace{0.4cm}

\noindent\textbf{2. 一个体育用品商店，每个排球 30 元，每个足球 50 元。李老师带的钱正好能买 12 个篮球或 10 个排球。（先根据问题填写表格，再列式解答）}

\vspace{0.6cm}

% 填写表格
\begin{center}
\renewcommand{\arraystretch}{1.8}
\begin{tabular}{|c|c|c|}
\noalign{\hrule height 1pt}
\makebox[3cm][c]{品名} & \makebox[4cm][c]{单价 / (元 / 个)} & \makebox[3cm][c]{数量 / 个} \\
\hline
排球 & \textcolor{blue}{\textbf{30}} & \textcolor{blue}{\textbf{10}} \\
\hline
篮球 & \textcolor{blue}{\textbf{?}} & \textcolor{blue}{\textbf{12}} \\
\hline
足球 & \textcolor{blue}{\textbf{50}} & \textcolor{blue}{\textbf{?}} \\
\noalign{\hrule height 1pt}
\end{tabular}
\end{center}

\vspace{0.8cm}

\noindent\textbf{（1）每个篮球多少元？}

\vspace{0.2cm}

{\small\color{gray}
\textbf{解析：}\\
第一步，先求出李老师带的总钱数：排球单价 $\times$ 排球数量 $= 30 \times 10 = 300$（元）。\\
第二步，由于这笔钱正好也能买 12 个篮球，所以篮球的单价为：总钱数 $\div$ 篮球数量 $= 300 \div 12 = 25$（元）。
}

\vspace{0.2cm}
\noindent\textbf{综合算式解答：} \textbf{\textcolor{blue}{$\boldsymbol{30 \times 10 \div 12 = 25}$}} \textbf{（元）} \\
\noindent\textbf{答：}每个篮球 \textbf{\textcolor{blue}{25}} 元。

\vspace{0.6cm}

\noindent\textbf{（2）李老师带的钱能买多少个足球？}

\vspace{0.2cm}

{\small\color{gray}
\textbf{解析：}\\
同样利用前面求出的总钱数 300 元。足球的单价是 50 元/个，那么这笔钱能买足球的数量为：总钱数 $\div$ 足球单价 $= 300 \div 50 = 6$（个）。
}

\vspace{0.2cm}
\noindent\textbf{综合算式解答：} \textbf{\textcolor{blue}{$\boldsymbol{30 \times 10 \div 50 = 6}$}} \textbf{（个）} \\
\noindent\textbf{答：}李老师带的钱能买 \textbf{\textcolor{blue}{6}} 个足球。

\vspace{0.6cm}

\noindent\textbf{（3）如果买 12 个篮球和 8 个足球，一共要多少元？}

\vspace{0.2cm}

{\small\color{gray}
\textbf{解析：}\\
买 12 个篮球所需的钱就是李老师原先的总钱数（$300$ 元），或者用刚才算出的篮球单价计算：$25 \times 12 = 300$（元）。\\
买 8 个足球需要的钱为：$50 \times 8 = 400$（元）。\\
两者加起来一共要：$300 + 400 = 700$（元）。
}

\vspace{0.2cm}
\noindent\textbf{综合算式解答：} \textbf{\textcolor{blue}{$\boldsymbol{25 \times 12 + 50 \times 8 = 700}$}} \textbf{（元）} \ \textit{或直接写为：} \textbf{\textcolor{blue}{$\boldsymbol{300 + 50 \times 8 = 700}$}} \textbf{（元）}\\
\noindent\textbf{答：}一共要 \textbf{\textcolor{blue}{700}} 元。


\newpage

\newpage

\subsection{解答：假设法解决问题——象棋与跳棋数量}

\vspace{0.4cm}

\noindent\textbf{3. 学校有象棋、跳棋共 26 副，2 人下一副象棋，6 人下一副跳棋，恰好可以供 96 人进行活动。象棋、跳棋各有多少副？（先假设两种棋的数量，再调整找出答案）}

\vspace{0.8cm}

\begin{center}
\renewcommand{\arraystretch}{1.8}
% 表格绘制：4列，包含假设过程、计算过程及结果对比
\begin{tabular}{|c|c|l|c|}
\hline
\textbf{象棋数量/副} & \textbf{跳棋数量/副} & \multicolumn{1}{c|}{\textbf{下棋总人数}} & \textbf{和 96 人比较} \\ \hline
\textcolor{blue}{13} & \textcolor{blue}{13} & \textcolor{blue}{$13 \times 2 + 13 \times 6 = 104$} & \textcolor{blue}{多 8 人} \\ \hline
\textcolor{blue}{14} & \textcolor{blue}{12} & \textcolor{blue}{$14 \times 2 + 12 \times 6 = 100$} & \textcolor{blue}{多 4 人} \\ \hline
\textcolor{blue}{15} & \textcolor{blue}{11} & \textcolor{blue}{$15 \times 2 + 11 \times 6 = 96$} & \textcolor{blue}{相等} \\ \hline
\end{tabular}
\end{center}

\vspace{0.6cm}

{\small\color{gray}
\textbf{解析：}\\
1. **第一次假设**：假设象棋和跳棋各有 $13$ 副（$26 \div 2 = 13$）。计算总人数为 $13 \times 2 + 13 \times 6 = 104$（人），比 $96$ 人多了 $8$ 人。\\
2. **分析调整方向**：因为计算出的总人数偏多，且跳棋的人均负荷（6人）高于象棋（2人），因此需要减少跳棋数量。\\
3. **逐步调整**：每减少一副跳棋并增加一副象棋，总人数减少 $6 - 2 = 4$（人）。\\
4. **得出结论**：经过两次调整，当象棋为 $15$ 副、跳棋为 $11$ 副时，总人数刚好为 $96$ 人。
}

\vspace{0.4cm}
\noindent\textbf{解答：}\\
象棋有（ \textbf{\textcolor{blue}{15}} ）副，跳棋有（ \textbf{\textcolor{blue}{11}} ）副。

\vspace{0.6cm}

{\small\color{blue!70!black}
\textbf{绘图基本原理与步骤：}\\
\textbf{1. 列表尝试法（Guess and Check）}：通过列表有序展示不同假设下的计算结果，寻找变量变化与最终数值间的规律。\\
\textbf{2. 建立逻辑关联}：表格中每一行代表一次完整的业务逻辑模拟。第一、二列为独立变量，第三列为衍生计算量，第四列为反馈修正量。\\
\textbf{3. 数据一致性}：表格内所有数据项均使用 \texttt{textcolor} 进行着色处理，既突出了动态求解过程，也保持了与题干风格的视觉统一。\\
\textbf{4. 结构化呈现}：利用 \texttt{tabular} 的水平线与垂直线构筑严谨的逻辑网格，使抽象的代数试错过程变得直观易读。
}



\newpage
